\documentclass[aspectratio=169]{beamer}
% Add slides directory to search path for Dartmouth theme
\makeatletter
\def\input@path{{../}}
\makeatother
\usetheme{Dartmouth}

% Packages
\usepackage{fontspec}
\usepackage{booktabs}
\usepackage{hyperref}
\usepackage{tikz}
\usetikzlibrary{shapes,arrows,positioning}

% Font settings
\setmonofont{Courier New}

% Title page information
\title{Theme Showcase}
\subtitle{Dartmouth Beamer Theme}
\author{PSYC 51.07}
\institute{Dr. Jeremy R. Manning}
\date{Template Deck}

\begin{document}

% ==============================================================================
% TITLE SLIDE
% ==============================================================================
\begin{frame}
\titlepage
\end{frame}

% ==============================================================================
% SECTION DIVIDER
% ==============================================================================
\section{Typography}

\begin{frame}[standout]
\centering
\Huge Typography
\end{frame}

% ==============================================================================
% SIMPLE BULLET SLIDE
% ==============================================================================
\begin{frame}{Simple Bullet Points}
\begin{itemize}
    \item First main point
    \item Second main point
    \item Third main point
\end{itemize}
\end{frame}

% ==============================================================================
% EMPHASIS EXAMPLES
% ==============================================================================
\begin{frame}{Text Emphasis}
\begin{itemize}
    \item Use \textbf{bold} for strong emphasis
    \item Use \emph{italics} for subtle emphasis
    \item Use \strong{Dartmouth green} for key terms
\end{itemize}
\end{frame}

% ==============================================================================
% ANIMATED BULLETS (opt-in example)
% ==============================================================================
\begin{frame}{Building Content Step-by-Step}
\begin{itemize}[<+->]
    \item This point appears first
    \item Then this point
    \item Finally this point
\end{itemize}

\onslide<4->
\begin{block}{Note}
Use animations sparingly for emphasis.
\end{block}
\end{frame}

% ==============================================================================
% BLOCK TYPES
% ==============================================================================
\section{Blocks}

\begin{frame}{Standard Block}
\begin{block}{Definition}
A language model predicts the probability of word sequences.
\end{block}
\end{frame}

\begin{frame}{Example Block}
\begin{exampleblock}{Example}
GPT-4 is a large language model trained on internet text.
\end{exampleblock}
\end{frame}

\begin{frame}{Alert Block}
\begin{alertblock}{Warning}
Always verify AI-generated content for accuracy.
\end{alertblock}
\end{frame}

% ==============================================================================
% TWO-COLUMN LAYOUT
% ==============================================================================
\section{Layouts}

\begin{frame}{Two-Column Layout}
\begin{columns}[T]
\begin{column}{0.45\textwidth}
\textbf{Left Column}
\begin{itemize}
    \item Point A
    \item Point B
\end{itemize}
\end{column}

\begin{column}{0.45\textwidth}
\textbf{Right Column}
\begin{itemize}
    \item Point X
    \item Point Y
\end{itemize}
\end{column}
\end{columns}
\end{frame}

% ==============================================================================
% TABLE
% ==============================================================================
\begin{frame}{Simple Table}
\begin{center}
\begin{tabular}{lcc}
\toprule
\textbf{Model} & \textbf{Year} & \textbf{Params} \\
\midrule
GPT-2 & 2019 & 1.5B \\
GPT-3 & 2020 & 175B \\
GPT-4 & 2023 & Unknown \\
\bottomrule
\end{tabular}
\end{center}
\end{frame}

% ==============================================================================
% CODE BLOCK
% ==============================================================================
\section{Code}

\begin{frame}[fragile]{Code Example}
\begin{dartmouthcode}[python]
def hello():
    """Simple function example."""
    print("Hello, Dartmouth!")
    return True
\end{dartmouthcode}
\end{frame}

% ==============================================================================
% EQUATION
% ==============================================================================
\section{Math}

\begin{frame}{Equation}
The attention mechanism:

\vspace{0.5cm}
\begin{center}
$\text{Attention}(Q, K, V) = \text{softmax}\left(\frac{QK^T}{\sqrt{d_k}}\right)V$
\end{center}
\end{frame}

% ==============================================================================
% SIMPLE TIKZ DIAGRAM
% ==============================================================================
\begin{frame}{Simple Diagram}
\begin{center}
\begin{tikzpicture}[
    node distance=2cm,
    box/.style={rectangle, draw=DartmouthGreen, fill=DartmouthGreen!10,
                minimum width=2cm, minimum height=0.8cm, thick}
]
\node[box] (input) {Input};
\node[box, right of=input] (process) {Process};
\node[box, right of=process] (output) {Output};

\draw[->, thick, DartmouthGreen] (input) -- (process);
\draw[->, thick, DartmouthGreen] (process) -- (output);
\end{tikzpicture}
\end{center}
\end{frame}

% ==============================================================================
% FINAL SLIDE
% ==============================================================================
\begin{frame}{Contact}
\centering
\Large Thank you!

\vspace{1cm}

\normalsize
jeremy@dartmouth.edu
\end{frame}

\end{document}
